\documentclass[a4paper,12pt]{article}
\usepackage[T1]{fontenc}
\usepackage[utf8]{inputenc}
\usepackage{amsmath}
\usepackage{amssymb}
\usepackage{enumitem}
\usepackage{amsthm}
\usepackage{geometry}
\usepackage{listings}
\usepackage{xcolor}

\geometry{margin=2cm}

\definecolor{codegray}{rgb}{0.5,0.5,0.5}
\definecolor{backcolour}{rgb}{0.95,0.95,0.92}

\lstdefinestyle{mystyle}{
    backgroundcolor=\color{backcolour},
    commentstyle=\color{codegray},
    keywordstyle=\color{blue},
    numberstyle=\tiny\color{codegray},
    stringstyle=\color{purple},
    basicstyle=\ttfamily\footnotesize,
    breakatwhitespace=false,
    breaklines=true,
    captionpos=b,
    keepspaces=true,
    showspaces=false,
    showstringspaces=false,
    showtabs=false,
    tabsize=4
}
\lstset{style=mystyle}

\begin{document}

\begin{center}
    \textbf{\Large PROGRAMOWANIE 1 R \hfill KOLOKWIUM \hfill 14.04.2025}
\end{center}

\vspace{0.5cm}
\noindent Name: [Imię, Nazwisko]

\vspace{0.5cm}

\noindent \textbf{1.} (2 punkty) Jaki będzie wynik wykonania poniższego kodu:
\begin{lstlisting}[language=Python]
t = [n for n in range(10)]
print(t[8:])
print(t[1:5:3])
print(t[5:1:-3])
\end{lstlisting}
\textcolor{blue}{\textit{Odpowiedź:}}
\begin{lstlisting}[language=Python, backgroundcolor=\color{white}]
[8, 9]
[1, 4]
[5, 2]
\end{lstlisting}

\vspace{0.5cm}

\noindent \textbf{2.} (4 punkty) Napisz jednolinijkowy, niewykorzystujący pętli ani instrukcji warunkowych kod wykonujący następujące zadania:

\noindent (A) wzajemna zamiana wartości zmiennych \texttt{x} i \texttt{y}, \\
\textcolor{blue}{\textit{Odpowiedź:}} \texttt{y, x = x, y}

\vspace{0.3cm}
\noindent (B) przypisanie do łańcucha tekstowego \texttt{t} wartości łańcucha (stringu) \texttt{s} zapisanej od tyłu, \\
\textcolor{blue}{\textit{Odpowiedź:}} \texttt{t = s[::-1]}

\vspace{0.3cm}
\noindent (C) przypisanie do zmiennej \texttt{d} listy dzielników liczby naturalnej przechowywanej w zmiennej \texttt{n}, \\
\textcolor{blue}{\textit{Odpowiedź:}} \texttt{d = [k for k in range(1, n+1) if n \% k == 0]}

\vspace{0.3cm}
\noindent (D) przypisanie do zmiennej \texttt{p} słownika, w którym klucze są liczbami naturalnymi podzielnymi przez 3 i mniejszymi od 25, zaś wartości - łańcuchami tekstowymi złożonymi z tylu liter \texttt{b}, ile wynosi wartość odpowiedniego klucza. \\
\textcolor{blue}{\textit{Odpowiedź:}} \texttt{p = \{k: (k * "b") for k in range(1, 26) if k \% 3 == 0\}}

\newpage
\begin{center}
    \textbf{\Large Programowanie 1 R \hfill Kolokwium \hfill Page 2 of 3}
\end{center}

\vspace{0.5cm}

\noindent \textbf{3.} (2 punkty) Obiektami których z wymienionych typów można indeksować słowniki? (można to nie znaczy że tak się powinno robić, tylko że interpreter nie zwróci błędu): \\
Zaznacz poprawne z odpowiedzi: \texttt{int}, \texttt{float}, \texttt{str}, \texttt{dict}, \texttt{tuple}, \texttt{function} \\
\textcolor{blue}{\textit{Odpowiedź (zaznaczone):}} \texttt{int}, \texttt{float}, \texttt{str}, \texttt{tuple}, \texttt{function}

\vspace{0.5cm}

\noindent \textbf{4.} (2 punkty) Co zostanie wypisane na ekranie w wyniku wykonania następującego kodu?
\begin{lstlisting}[language=Python]
x = ["a", "b"]
y = x
y[0], y[1] = x[1], x[0]
print(x+y)
\end{lstlisting}
\textcolor{blue}{\textit{Odpowiedź:}} \texttt{baba} \textit{(zapisane odręcznie w ten sposób, wynik z interpretera to \texttt{['b', 'a', 'b', 'a']})}

\vspace{0.5cm}

\noindent \textbf{5.} (2 punkty) Co zostanie wypisane na ekranie w wyniku wykonania następującego kodu?
\begin{lstlisting}[language=Python]
listA = list(filter(lambda x: (x%2 != 0), [c//2 for c in range(20)]))
listB = list(map(lambda x: x * 2, listA))
print(listB)
\end{lstlisting}
\textcolor{blue}{\textit{Odpowiedź:}} \texttt{2 2 6 6 10 10 14 14 18 18} \textit{(zapisane odręcznie w ten sposób, wynik z interpretera to \texttt{[2, 2, 6, 6, 10, 10, 14, 14, 18, 18]})}

\vspace{0.5cm}

\noindent \textbf{6.} (2 punkt) Jeśli dana jest lista \texttt{a} zawierająca elementy typu \texttt{str} to użycie metody \texttt{a.sort()} posortuje listę \texttt{a} w porządku "słownikowym". Jak użyć metody \texttt{a.sort()} aby posortowała ona listę w porządku rosnącej liczby liter? \\
\textcolor{blue}{\textit{Odpowiedź:}} \texttt{a.sort(key=len)}

\newpage
\begin{center}
    \textbf{\Large Programowanie 1 R \hfill Kolokwium \hfill Page 3 of 3}
\end{center}

\vspace{0.5cm}

\noindent \textbf{7.} (2 punkty) Wskaż błąd w poniższym programie, zakładając że wynikiem ma być 549:
\begin{lstlisting}[language=Python]
data = [{'id': 1, 'data': 'jeden'},
        {'id': 33, 'data': 'trzydziesci trzy'},
        {'id': 515, 'data': 'piecset pietnascie'}]
f = lambda x: x['id']
l = [f for f in data]
print(sum(l))
\end{lstlisting}
\textcolor{blue}{\textit{Odpowiedź (poprawka odręczna):}} Zamiast \texttt{l = [f for f in data]} powinno być \texttt{l = [f(x) for x in data]}

\vspace{0.5cm}

\noindent \textbf{8.} (2 punkty) Dana jest zmienna \texttt{s} typu \texttt{str} (która składa się tylko z liter tworzących słowo). Napisz fragment kodu tworzący słownik, którego klucze to litery występujące w tym słowie, a odpowiadające im wartości to liczby powtórzeń tych liter. \\
\textcolor{blue}{\textit{Odpowiedź:}} \texttt{\{c: s.count(c) for c in set(list(s))\}}

\vspace{0.5cm}

\noindent \textbf{9.} (2 punkty) Co zostanie wypisane na ekranie po wykonaniu poniższego kodu?
\begin{lstlisting}[language=Python]
def mystery(n):
    return [i * j for i in range(1, n) for j in range(i, n) if (i + j) % 3 == 0]
print(mystery(5))
\end{lstlisting}
\textcolor{blue}{\textit{Odpowiedź:}} \texttt{[2, 8, 9]} \textit{(odręcznie przekreślone 12 i 15 na końcu)}

\end{document}